\chapter{Fazit und Ausblick}
\label{ch:fazit}

% ==============================================================================
\section{Fazit}
\label{sec:fazit}

\todo[inline]{Ergebnis steht noch aus. Nach den Versuchen mit dem neuen Datensatz schreiben:
Erfolgsquote an der Anlage, Genauigkeit, Vergleich mit dem klassischen Greifen von Modul~2.}

Bis zum Stand dieses Dokuments ist die vollständige Kette aus automatischer Aufnahme mit
eingespeistem Rauschen, Training einer Diffusion Policy und sicherer Ausführung umgesetzt und
in das Gesamtsystem eingebunden. An der Anlage findet die Policy das Bauteil an unbekannten
Lagen aus den Bildern und trifft den Greifpunkt im Median auf knapp \SI{1}{\centi\metre}.
Die verbleibenden Fehler entstehen im Endanflug und bei der Drehung, in Zuständen, die die
Demonstrationen nicht abdecken.

% ==============================================================================
\section{Ausblick}
\label{sec:ausblick}

\paragraph{DAgger.} Die Fehler an der Anlage entstehen dort, wo die Demonstrationen keine
Beispiele haben. Genau diese Zustände sammelt \ac{dagger} (\cref{sec:covariate-shift}). Da der
Experte in Modul~4 ein Planer ist, kann er jeden Zustand nachträglich beschriften: Bei
bekannter Lage ist die richtige Aktion \enquote{von hier senkrecht über den Greifpunkt, dann
ab}. Geplant ist, Fahrten der Policy mit Bildern aufzuzeichnen, offline mit den Aktionen des
Planers zu beschriften und im Training höher zu gewichten. Erste Episoden dieser Art wurden am
17.11.\ aufgenommen \quelle{Dokumentation/DAgger-Ablauf.md}.

\paragraph{Weitere Punkte.}
\begin{itemize}
  \item Gegenversuch mit dem gewichteten Modell C3 an den Minus-Lagen,
  \item Drehkorrektur unterhalb von \code{ANFLUG} lehren und die Drehfahrt verlangsamen,
  \item geordneter Abbruch aus dem Ablauf des Gesamtsystems,
  \item Vergleich nur mit Handgelenkkamera, da der Datensatz beide Bildströme enthält,
  \item Übertragung auf weitere Bauteile und auf das Greifen aus einer Kiste.
\end{itemize}
